	A crescente adoção de ferramentas de inteligência artificial generativa 
baseadas em modelos de linguagem de grande escala (LLMs) tem transformado o 
desenvolvimento de software, deslocando o esforço do desenvolvedor da escrita do 
código para as etapas de revisão e auditoria. Este trabalho avaliuou 
qualitativamente o código-fonte gerado por LLMs. Adota-se uma abordagem 
qualitativa, fundamentada em revisão bibliográfica e em testes empíricos que 
comparam código gerado por inteligência artificial e código escrito manualmente. 

\textbf{Palavras-chave}: Inteligência artificial; modelos de linguagem de grande 
escala; qualidade de código; viés de automação; engenharia de software. 